\documentclass[a4paper,11pt]{article}
\usepackage[utf8]{inputenc}
\usepackage[vietnamese]{babel}
\usepackage{amsmath,amssymb}
\usepackage{mathptmx}
\usepackage[margin=2.2cm]{geometry}
\usepackage{booktabs}
\usepackage{tabularx}
\usepackage{ragged2e}
\usepackage{enumitem}
\usepackage{microtype}
\usepackage{listings}
\usepackage{xcolor}
\usepackage[hidelinks]{hyperref}

\newcolumntype{Y}{>{\RaggedRight\arraybackslash}X}
\newcolumntype{L}[1]{>{\RaggedRight\arraybackslash}p{#1}}

\lstset{
  basicstyle=\ttfamily\small,
  breaklines=true,
  frame=single,
  backgroundcolor=\color{gray!5},
  rulecolor=\color{gray!40},
  keywordstyle=\bfseries\color{blue!70!black},
  commentstyle=\itshape\color{green!50!black},
  stringstyle=\color{red!60!black},
  showstringspaces=false
}

\title{\textbf{Chương 3: Giải Phẫu Tech Stack Thực Chiến}\\[0.3em]
\large\textit{Làm Chủ Playwright, Công Cụ Thay Thế Tối Ưu, Sức Mạnh Của Python và Vòng Lặp Tự Sửa Lỗi}}
\author{\textbf{Giáo Trình Kỹ Nghệ Vibecoding Chuyên Sâu 2026}}
\date{Tháng 10 Năm 2026}

\begin{document}

\maketitle

\begin{abstract}
Khác với các tài liệu lý thuyết bề mặt, Chương 3 là bản giải phẫu thực hành chuyên sâu về các công cụ cốt lõi trong quy trình Vibecoding tác tử tự chủ (\textit{Agentic Loop}). Chương học mổ xẻ 7 siêu năng lực thực chiến của Playwright, phân tích các điểm nghẽn bộ nhớ, đồng thời so sánh trực diện với các công cụ tối ưu hơn như Crawl4AI, Browser-use và HTTPX. Tiếp đó, tài liệu cung cấp mã nguồn chuẩn cho 2 luồng công việc cốt tử: tự động thu thập bài báo nghiên cứu và tự động kiểm thử webapp kèm cơ chế xuất phản hồi máy đọc (\texttt{agent\_feedback.json}) để AI tự sửa lỗi; đồng thời xác lập vị thế của Python như một chất keo tự động hóa tối thượng trong tương quan với TypeScript.
\end{abstract}

\vspace{1em}
\rule{\textwidth}{0.4pt}

\section{Playwright Thực Chiến: Khả Năng và Bảng Đối Chiếu Công Cụ Tối Ưu}

Trong quy trình Vibecoding hiện đại, Playwright đóng vai trò là ``đôi mắt và bàn tay'' giúp mô hình AI tương tác với giao diện thực tế.

\subsection{Bảy siêu năng lực thực tế của Playwright}
\begin{enumerate}[label={\arabic*.}]
    \item \textbf{Headless Browser Automation}: Khởi chạy đa nhân Chromium, Firefox, WebKit ngầm không cần giao diện, thích hợp tuyệt đối trên máy chủ Linux/CI.
    \item \textbf{Lắng nghe \& Bắt chặn gói tin (Network Interception)}: Bắt trộm dữ liệu máy khách gửi lên (\texttt{page.on('request')}) và giả lập phản hồi API (\texttt{page.route()}) để kiểm thử tình huống máy chủ bị lỗi 500 hoặc mạng chập chờn.
    \item \textbf{Cơ chế tự động chờ (Auto-waiting)}: Tự động chờ phần tử hiển thị, sẵn sàng nhận thao tác mà không cần dùng lệnh chờ cứng \texttt{time.sleep()}.
    \item \textbf{Bắt lỗi Console \& Ngoại lệ chưa xử lý}: Bắt toàn bộ \texttt{console.error} và \texttt{pageerror} (như React Hydration mismatch, Null pointer crash).
    \item \textbf{Chụp ảnh màn hình kiểm thử thị giác (Visual Diff)}: Chụp ảnh full-page làm minh chứng đối chiếu giao diện cho AI Agent.
    \item \textbf{Lưu trữ phiên đăng nhập (Storage State)}: Ghi nhận cookies và localStorage vào tệp \texttt{auth.json} sau một lần đăng nhập để tái sử dụng.
    \item \textbf{Trích xuất DOM động từ Single Page Applications (SPA)}: Chờ ứng dụng React/Next.js hoàn tất render dữ liệu máy khách rồi mới trích xuất nội dung.
\end{enumerate}

\subsection{Bảng so sánh: Khi nào Playwright quá tải và Các công cụ thay thế}
Mỗi phiên Chromium ngốn từ 250MB đến 600MB RAM. Việc lạm dụng Playwright cho mọi tác vụ đọc web là một sai lầm nghiêm trọng gây tràn RAM hệ thống (\textit{OOM}).

\begin{table}[htbp]
\centering
\small
\caption{So sánh Playwright với các công cụ thu thập và tự động hóa thay thế}
\vspace{0.5em}
\begin{tabularx}{\textwidth}{L{2.5cm} L{2cm} L{2.2cm} L{3.5cm} Y}
\toprule
\textbf{Công cụ} & \textbf{Tiêu thụ RAM} & \textbf{Tốc độ} & \textbf{Đặc tính cốt lõi} & \textbf{Khi nào nên sử dụng?} \\
\midrule
\textbf{Playwright} & 300--600 MB & 2--5 giây & Mô phỏng đầy đủ Chromium, click DOM, bắt console log & Bắt buộc khi test WebApp tương tác, bắt lỗi console, chụp ảnh giao diện \\
\textbf{Crawl4AI} & 80--150 MB & 0.5--1.5 giây & Chuyên cho LLM, tự chuyển HTML thành Markdown sạch & Rất hay khi cào bài báo, blog, tài liệu để nhồi vào prompt của AI \\
\textbf{Browser-Use / Stagehand} & 400--700 MB & Chậm (gọi AI) & Tác tử tự nhìn màn hình và tự click bằng ngôn ngữ tự nhiên & Dùng khi web có cấu trúc quá phức tạp, selector CSS thay đổi liên tục \\
\textbf{HTTPX + BeautifulSoup} & \textbf{10--20 MB} & \textbf{$< 0.2$ giây} & Client HTTP thuần túy, không tốn tài nguyên trình duyệt & Rất tối ưu khi cào arXiv, Wikipedia, API tĩnh; tiết kiệm 95\% RAM \\
\bottomrule
\end{tabularx}
\end{table}

\section{Hai Quy Trình Thực Chiến Bằng Mã Nguồn Mẫu}

\subsection{Quy trình 1: Thu thập và Thẩm định Bài báo Nghiên cứu Tự động}
Thay vì mở trình duyệt nặng nề, kỹ sư sử dụng HTTPX gọi thẳng arXiv API để lấy cấu trúc dữ liệu XML chỉ trong 0.8 giây; chỉ khi cần chụp ảnh minh chứng thị giác mới khởi chạy Playwright ngầm:

\begin{lstlisting}[language=python, caption={Kich ban thu thap bai bao sieu nhe ket hop Playwright va HTTPX}]
import httpx, xml.etree.ElementTree as ET
from playwright.sync_api import sync_playwright

def fetch_arxiv_fast(paper_id: str):
    url = f"https://export.arxiv.org/api/query?id_list={paper_id}"
    resp = httpx.get(url, follow_redirects=True, timeout=10.0)
    root = ET.fromstring(resp.text)
    entry = root.find('{http://www.w3.org/2005/Atom}entry')
    title = entry.find('{http://www.w3.org/2005/Atom}title').text.strip()
    abstract = entry.find('{http://www.w3.org/2005/Atom}summary').text.strip()
    return {"id": paper_id, "title": title, "abstract": abstract}

def verify_paper_screenshot(paper_id: str, out_img: str):
    with sync_playwright() as p:
        browser = p.chromium.launch(headless=True)
        page = browser.new_page()
        page.goto(f"https://arxiv.org/abs/{paper_id}", timeout=15000)
        page.screenshot(path=out_img)
        browser.close()
\end{lstlisting}

\subsection{Quy trình 2: Tự Động Test WebApp \& Vòng Lặp Tự Sửa Lỗi (Self-Healing)}
Kịch bản lắng nghe toàn bộ 4 kênh lỗi: console.error, uncaught pageerror, network 4xx/5xx và requestfailed. Khi phát hiện lỗi, script ghi dữ liệu vào tệp máy đọc \texttt{agent\_feedback.json} kèm số dòng cụ thể và thoát với Exit Code 1. AI Agent đọc tệp này để tự khoanh vùng và sửa code cho đến khi Exit Code là 0:

\begin{lstlisting}[language=python, caption={Co che xuat phan hoi loi co cau truc cho AI Agent}]
# scripts/self_healing_test_runner.py (Core Implementation)
errors = []
page.on("console", lambda m: errors.append({"type": "console", "text": m.text, "loc": m.location}) if m.type == "error" else None)
page.on("pageerror", lambda e: errors.append({"type": "exception", "msg": str(e)}))
page.on("response", lambda r: errors.append({"type": "http", "status": r.status, "url": r.url}) if r.status >= 400 else None)

page.goto("http://localhost:3000", wait_until="networkidle")
if errors:
    page.screenshot(path="failure_screenshot.png")
    with open("agent_feedback.json", "w") as f:
        json.dump({"status": "FAILED", "errors": errors}, f, indent=2)
    sys.exit(1) # Exit Code 1: Agent continues fixing code
sys.exit(0)     # Exit Code 0: All green, loop completed!
\end{lstlisting}

\section{Sức Mạnh Toàn Diện Của Python Trong Kỷ Nguyên Vibecoding}

Python không chỉ là ngôn ngữ khoa học dữ liệu, mà trong Vibecoding nó chính là \textbf{chất keo tự động hóa tối thượng} với 6 trụ cột:
\begin{enumerate}[label={\arabic*.}]
    \item \textbf{Tự động hóa kịch bản \& File system}: Quản lý tiến trình, biên dịch tài liệu LaTeX tự động, dọn dẹp các tiến trình Chrome ma bị treo ngầm.
    \item \textbf{Cào dữ liệu \& Xử lý tri thức}: Tích hợp nhịp nhàng giữa HTTPX (nhẹ, nhanh) và Playwright (vượt JS) để nạp tài liệu cho AI.
    \item \textbf{Backend hiệu năng cao với FastAPI \& Pydantic}: Xây dựng API bất đồng bộ siêu tốc, tự động sinh tài liệu Swagger/OpenAPI và kiểm duyệt dữ liệu nghiêm ngặt.
    \item \textbf{Xử lý PDF \& Tài liệu chuyên sâu}: Bóc tách từng trang PDF đề thi bằng \texttt{pypdf}, \texttt{pdfplumber}, xuất công thức Toán KaTeX/LaTeX chuẩn mực.
    \item \textbf{Xử lý dữ liệu cấu trúc}: Dùng Polars và Pandas biến đổi hàng nghìn dòng dữ liệu JSON/CSV chỉ trong một phần nghìn giây.
    \item \textbf{Điều phối AI cục bộ \& Embeddings}: Kết nối các mô hình cục bộ qua LiteLLM, llama.cpp và tính toán vector ngữ nghĩa.
\end{enumerate}

\begin{table}[htbp]
\centering
\small
\caption{Bảng so sánh chiến lược: Khi nào dùng Python vs TypeScript khi Vibecode}
\vspace{0.5em}
\begin{tabularx}{\textwidth}{L{3.5cm} L{3.5cm} Y}
\toprule
\textbf{Lĩnh vực / Tác vụ} & \textbf{Ngôn ngữ tối ưu} & \textbf{Lý do kỹ thuật} \\
\midrule
\textbf{Frontend Web/Mobile UI} & \textbf{TypeScript} (React/Next) & Hệ sinh thái component cực mạnh (Tailwind, shadcn/ui), JSX/TSX trực quan \\
\textbf{Kịch bản tự động hóa \& Agent Tools} & \textbf{Python} & Viết 30 dòng chạy ngay, không cần cấu hình \texttt{tsconfig.json} hay bước build \\
\textbf{Dịch vụ AI Backend, Toán, PDF} & \textbf{Python} (FastAPI) & Toàn bộ thư viện AI/ML và trích xuất PDF/LaTeX đỉnh cao nằm ở Python \\
\textbf{Ứng dụng Full-Stack Monorepo} & \textbf{TypeScript} (tRPC) & Đảm bảo End-to-End Type Safety xuyên suốt từ database đến nút bấm UI \\
\bottomrule
\end{tabularx}
\end{table}

\section{Tầng Giao Tiếp API và Lưu Trữ: Ranh Giới Kỹ Thuật}

\subsection{Cạm bẫy CORS và Serverless Timeout khi gọi AI}
\begin{itemize}[leftmargin=1.5em]
    \item \textbf{CORS Misconfiguration}: Không được phép gán bừa \texttt{Access-Control-Allow-Origin: *} kèm credentials vì sẽ mở toang lỗ hổng cho tấn công CSRF.
    \item \textbf{Serverless Timeout (504)}: Khi gọi LLM phân tích văn bản mất hơn 15 giây, môi trường Serverless (như Vercel) sẽ tự ngắt kết nối. Bắt buộc phải áp dụng kỹ thuật truyền phát trực tiếp (\textit{Server-Sent Events - SSE Streaming}) hoặc hàng đợi nền (\textit{Background Queues}).
\end{itemize}

\subsection{Supabase PostgreSQL đối đầu SQLite}
Supabase mang lại sự an tâm cho các ứng dụng đa người dùng nhờ cơ chế phân quyền bảo mật mức hàng (\textit{Row Level Security - RLS}) nhưng đòi hỏi quản lý Connection Pooling qua PgBouncer/Supavisor. Ngược lại, SQLite cực kỳ hoàn hảo cho các bản MVP chạy cục bộ siêu nhẹ nhưng sẽ bị nghẽn khóa ghi (\textit{File lock}) khi có nhiều luồng ghi cùng lúc.

\subsection{Cây quyết định: Lựa chọn Công cụ Thu thập Dữ liệu và Lưu trữ}

Người học không nên mặc định dùng Playwright cho mọi trường hợp cào dữ liệu hay Supabase cho mọi bài toán lưu trữ. Dưới đây là ma trận quyết định giúp tối ưu tài nguyên và tốc độ thực thi:

\begin{table}[htbp]
\centering
\small
\caption{Cây quyết định lựa chọn công cụ thu thập dữ liệu \& lưu trữ}
\vspace{0.5em}
\begin{tabularx}{\textwidth}{L{2.8cm} L{3.5cm} L{4.5cm} Y}
\toprule
\textbf{Nhu cầu bài toán} & \textbf{Công cụ khuyến nghị} & \textbf{Cơ chế hoạt động tối ưu} & \textbf{Hạn chế kỹ thuật} \\
\midrule
\textbf{Lấy dữ liệu API / Trang tĩnh} & \textbf{HTTPX} (Async) & Gửi HTTP GET/POST trực tiếp, chỉ tốn 15MB RAM, tốc độ $< 0.8$s & Bất lực trước SPA render bằng JS (React/Vue/Angular) \\
\textbf{Bóc tách trang web thành Markdown} & \textbf{Jina Reader} (\texttt{r.jina.ai}) & Thêm tiền tố \texttt{r.jina.ai/URL}, trả về Markdown sạch 0 dòng code & Phụ thuộc dịch vụ bên ngoài, giới hạn rate limit \\
\textbf{Vượt Bot Detection / Cloudflare} & \textbf{Nodriver} / \textbf{Crawl4AI} & Điều khiển Chrome DevTools Protocol (CDP) không gắn cờ webdriver & Cấu hình phức tạp hơn Playwright truyền thống \\
\textbf{Test Webapp \& Chụp màn hình} & \textbf{Playwright} & Bắt toàn diện 4 kênh lỗi (Console, PageError, Network 4xx/5xx) & Tốn 300--500MB RAM cho mỗi tiến trình trình duyệt \\
\midrule
\textbf{Lưu trữ Cục bộ / Đơn người dùng} & \textbf{SQLite} (WAL mode) & 1 tệp tin duy nhất, tốc độ đọc triệu dòng/s, tích hợp FTS5 & Nghẽn ghi khi có nhiều người dùng ghi đồng thời \\
\textbf{Phân tích Dữ liệu / OLAP} & \textbf{DuckDB} & Tối ưu truy vấn cột (\textit{columnar}), đọc trực tiếp Parquet/CSV & Không phù hợp cho ứng dụng web giao dịch (OLTP) \\
\textbf{Backend Tức thì / Auth + DB} & \textbf{PocketBase} & 1 file nhúng nhị phân duy nhất, có sẵn Auth, Storage và Realtime & Không mở rộng quy mô đa máy chủ ngang (horizontal scale) \\
\textbf{Production SaaS / Đa người dùng} & \textbf{PostgreSQL} (Supabase) & Đầy đủ RLS, ACID, PgBouncer, quan hệ khóa ngoại phức tạp & Cần cấu hình kết nối cẩn trọng để tránh cạn Connection Pool \\
\bottomrule
\end{tabularx}
\end{table}

\subsection{Kịch bản thực hành đồng hành trong thư mục}
Trong thư mục chương 3 có sẵn 2 công cụ thực chiến hoàn chỉnh:
\begin{itemize}[leftmargin=1.5em]
    \item \texttt{paper\_harvester.py}: Công cụ kép thu hoạch bài báo khoa học arXiv (HTTPX siêu nhẹ + Playwright dự phòng).
    \item \texttt{self\_healing\_test\_runner.py}: Bộ chạy kiểm thử Playwright tự động quét 4 kênh lỗi và xuất \texttt{agent\_feedback.json} phục vụ vòng lặp tự sửa lỗi.
\end{itemize}

\section{Tổng kết chương và bài tập thực hành}

\subsection{Tóm tắt cốt lõi}
\begin{itemize}[leftmargin=1.5em]
    \item Không lạm dụng Playwright; luôn ưu tiên HTTPX API khi có thể để tiết kiệm 95\% RAM.
    \item Sử dụng Playwright có mục đích: bắt toàn bộ 4 kênh lỗi và xuất ra tệp \texttt{agent\_feedback.json} để tạo vòng lặp tự sửa lỗi khép kín.
    \item Tận dụng Python làm chất keo tự động hóa kịch bản, xử lý PDF/AI backend; kết hợp cùng TypeScript cho tầng giao diện.
\end{itemize}

\subsection{Bài tập thực hành}
\begin{enumerate}[label={\arabic*.}]
    \item \textbf{Chạy Paper Harvester}: Sử dụng \texttt{scripts/paper\_harvester.py} để tìm kiếm 3 bài báo mới nhất về ``Autonomous Agents'' và tải về thư mục cục bộ.
    \item \textbf{Mô phỏng Self-Healing}: Chạy \texttt{scripts/self\_healing\_test\_runner.py} trên một trang web bất kỳ của bạn, quan sát cách tệp \texttt{agent\_feedback.json} được tạo ra và kiểm tra mã lỗi Exit Code.
\end{enumerate}

\end{document}
